\chapter{Contractions}

\section*{Definition and Overview}
Uterine contractions are the periodic, involuntary tightening and shortening of the uterine myometrium. Throughout pregnancy, contractions serve various physiological purposes: during gestation, irregular and mild contractions (known as Braxton Hicks contractions) help prepare the uterus for birth. During labour, coordinated, regular, and progressively stronger contractions are the driving force behind cervical effacement and dilation, facilitating the descent and delivery of the fetus, the expulsion of the placenta, and subsequent uterine involution to prevent postpartum haemorrhage. The clinical significance of monitoring contractions lies in evaluating the progress of labour, detecting preterm labour, and assessing fetal wellbeing through its relationship with the fetal heart rate.

\section*{Epidemiology}
Uterine contractions are a universal physiological experience for all women who reach a viable stage of pregnancy. Braxton Hicks contractions are commonly felt from the second trimester onwards, reported by approximately 75\% of pregnant individuals. Active labour contractions occur at term (between 37 and 42 weeks of gestation) in the vast majority of pregnancies. Preterm contractions leading to preterm birth affect approximately 8\% of births in Australia. Monitoring and managing abnormal contraction patterns (such as hyperstimulation or incoordinate labour) is a daily component of obstetric and midwifery care in birth units globally.

\section*{Aetiology and Pathophysiology}
The initiation and propagation of uterine contractions are governed by a complex cascade of endocrine, mechanical, and cellular events.
\begin{itemize}
    \item \textbf{Oxytocin receptor upregulation:} Under the influence of rising oestrogen levels, the density of oxytocin receptors in the myometrium increases dramatically towards the end of pregnancy, making the uterus highly sensitive to pulsatile oxytocin released from the posterior pituitary.
    \item \textbf{Progesterone withdrawal:} A local shift in the progesterone-to-oestrogen ratio decreases the relaxing effect of progesterone on the myometrium, promoting uterine excitability.
    \item \textbf{Prostaglandin synthesis:} The decidua and fetal membranes produce prostaglandins (specifically PGE2 and PGF2$\alpha$), which stimulate myometrial contractions and promote cervical ripening.
    \item \textbf{Myometrial gap junction formation:} Estrogen stimulates the expression of connexin 43, forming gap junctions between myometrial cells. This allows rapid electrical coupling and synchronized, coordinated contractions across the entire uterus.
    \item \textbf{Calcium ion influx:} Depolarisation of myometrial cells opens voltage-gated calcium channels, leading to an influx of calcium that activates myosin light-chain kinase, causing actin-myosin cross-bridge cycling and muscle contraction.
\end{itemize}

\section*{Clinical Features}
True labour contractions are characterized by their regular frequency, progressive intensity, and effectiveness in changing the cervix, contrasting sharply with false labour.
\begin{itemize}
    \item \textbf{Regular frequency:} Contractions occur at regular intervals, progressively shortening in time between each episode (e.g., moving from 10 minutes apart to 2 to 3 minutes apart).
    \item \textbf{Increasing duration:} Each contraction lasts between 30 and 70 seconds, gradually lengthening as labour advances.
    \item \textbf{Progressive intensity:} The physical pain and tightening become stronger over time, starting as a dull ache in the lower back and radiating around to the lower abdomen. They do not subside with rest, hydration, or position changes.
    \item \textbf{Associated cervical changes:} Accompanied by cervical effacement (thinning) and dilation, often accompanied by the loss of the mucous plug (the "show").
    \item \textbf{Braxton Hicks features:} Typically irregular, infrequent (fewer than 4 per hour), short in duration, painless or mildly uncomfortable, felt only in the front of the abdomen, and resolve with rest, walking, or drinking water.
\end{itemize}

\section*{Differential Diagnosis}
Distinguishing true labour contractions from other causes of abdominal pain or cramping is critical for appropriate maternal and fetal management.
\begin{itemize}
    \item \textbf{Braxton Hicks contractions:} Non-productive, irregular uterine contractions that do not result in cervical change.
    \textbf{Preterm labour:} Regular uterine contractions causing cervical change occurring prior to 37 completed weeks of gestation.
    \item \textbf{Placental abruption:} Characterised by constant, severe abdominal pain, uterine tenderness or rigidity (hypertonus), and often vaginal bleeding, rather than intermittent contraction pain.
    \item \textbf{Urinary tract infection (UTI) or pyelonephritis:} Can cause lower abdominal cramping, pelvic discomfort, back pain, dysuria, and systemic symptoms like fever.
    \item \textbf{Gastrointestinal pathology:} Conditions such as appendicitis, gastroenteritis, irritable bowel syndrome, or constipation can mimic or trigger uterine cramping.
\end{itemize}

\section*{When to Seek Medical Care}
Pregnant women should contact their midwife, obstetrician, or maternity hospital if they experience regular contractions, vaginal bleeding, a gush or trickle of fluid (rupture of membranes), or a reduction in fetal movements.

\textbf{Call 000 (or local emergency services) for:}
\begin{itemize}
    \item Sudden, heavy or bright red vaginal bleeding
    \item Constant, severe abdominal pain without relief between contractions
    \item Signs of umbilical cord prolapse (feeling the cord at the vaginal opening after water breaks)
    \item Severe headache, sudden swelling of the face, hands, or feet, or visual disturbances (blurring, flashing lights)
    \item Maternal collapse, difficulty breathing, or seizures
\end{itemize}

\section*{Diagnosis and Investigations}
Evaluating contractions involves confirming labour, monitoring fetal tolerance, and ruling out complications.
\begin{itemize}
    \item \textbf{Clinical palpation:} Placing a hand on the maternal abdomen (fundus) to assess the frequency, duration, and strength (mild, moderate, or strong/board-like) of contractions.
    \item \textbf{Cardiotocography (CTG):} Continuous electronic fetal monitoring that simultaneously records the fetal heart rate and uterine contractions (via a tocodynamometer), helping identify contraction patterns and fetal distress.
    \item \textbf{Digital cervical examination:} Assessment of the cervix to determine dilation (in centimetres), effacement (percentage), position (posterior, mid, anterior), and consistency (soft or firm).
    \item \textbf{Fetal fibronectin (fFN) test:} A vaginal swab test used in symptomatic women between 22 and 34 weeks of gestation to assess the likelihood of preterm delivery within the next 7 to 14 days.
    \item \textbf{Transvaginal ultrasound:} Used to measure cervical length in cases of suspected preterm labour; a shortened cervix increases the risk of preterm delivery.
\end{itemize}

\section*{Management}
The management of contractions depends on the gestational age and whether the goal is to facilitate or suppress uterine activity.
\begin{itemize}
    \item \textbf{Non-pharmacological pain relief (Active Labour):} Utilising warm baths (hydrotherapy), massage, heat packs, active birth positions (upright, kneeling, leaning forward), TENS machines, and breathing techniques.
    \item \textbf{Pharmacological analgesia (Active Labour):} Nitrous oxide gas (gas and air), intramuscular or intravenous opioids (e.g., morphine, fentanyl), or epidural anaesthesia for complete pain block.
    \item \textbf{Tocolytic therapy (Preterm Labour):} Administration of medications such as nifedipine (calcium channel blocker) or atosiban (oxytocin receptor antagonist) to temporarily suppress contractions, allowing time for corticosteroid administration to mature fetal lungs and transfer to a tertiary facility.
    \item \textbf{Labour induction and augmentation:} The use of synthetic oxytocin (Syntocinon) infusion or amniotomy (artificial rupture of membranes) to stimulate or strengthen contractions when labour is prolonged or induced.
    \item \textbf{Active management of third stage:} Administration of an uterotonic (e.g., oxytocin) immediately after the birth of the baby to stimulate strong contractions, facilitating placental delivery and preventing postpartum haemorrhage.
\end{itemize}

\section*{Complications}
Abnormal contraction patterns can lead to maternal and fetal complications during pregnancy and childbirth.
\begin{itemize}
    \item \textbf{Uterine hyperstimulation/tachysystole:} Defined as more than 5 contractions in a 10-minute period. This can restrict uteroplacental blood flow, leading to fetal hypoxia and distress.
    \item \textbf{Incoordinate uterine action:} Contractions that are irregular, painful, but ineffective in dilating the cervix, leading to prolonged labour and maternal exhaustion.
    \item \textbf{Preterm delivery:} Contractions that occur before 37 weeks, resulting in neonatal prematurity and associated respiratory, neurological, and developmental risks.
    \item \textbf{Uterine rupture:} A rare but catastrophic tearing of the uterine wall, usually associated with a prior caesarean section scar and augmented contractions, leading to severe maternal and fetal haemorrhage.
    \item \textbf{Uterine atony:} Failure of the uterus to contract effectively after birth, which is the most common cause of major postpartum haemorrhage.
\end{itemize}

\section*{Prognosis and Outcomes}
For the majority of term pregnancies, uterine contractions progress naturally to a successful vaginal delivery. The active first stage of labour averages 8 to 18 hours for a first pregnancy and 5 to 12 hours for subsequent deliveries. Following birth, contractions continue in a mild, irregular form (known as afterbirth pains) for several days, especially during breastfeeding due to endogenous oxytocin release, aiding the uterus in returning to its pre-pregnancy size. In cases of preterm labour, prompt administration of tocolytics and corticosteroids significantly improves neonatal survival and reduces respiratory distress syndrome.

\section*{Prevention and Self-Care}
Self-care measures focus on managing discomfort and preventing premature uterine activity.
\begin{itemize}
    \item \textbf{Maintain adequate hydration:} Dehydration is a common trigger for Braxton Hicks contractions; drinking plenty of water throughout the day can calm uterine irritability.
    \item \textbf{Practice position changes and rest:} If experiencing Braxton Hicks, lying down or changing positions (such as taking a gentle walk) often resolves the tightening.
    \item \textbf{Attend childbirth preparation classes:} Education on labor stages, breathing, and partner support helps reduce anxiety, which can otherwise impede natural oxytocin release.
    \item \textbf{Use heat therapy:} Applying a warm heat pack to the lower back or taking a warm bath can soothe mild cramping.
    \item \textbf{Avoid heavy lifting or overexertion:} Excessive physical strain in late pregnancy can trigger frequent Braxton Hicks or premature contractions.
\end{itemize}

\section*{Sources and Further Reading}
The following reputable organisations provide further information on uterine contractions and labour:
\begin{itemize}
    \item Pregnancy, Birth and Baby (Australian Government). \textit{Understanding contractions and labour.} https://www.pregnancybirthbaby.org.au/
    \item Royal Australian and New Zealand College of Obstetricians and Gynaecologists (RANZCOG). \textit{Labour and birth resources.} https://ranzcog.edu.au/
    \item Queensland Health. \textit{Maternity care clinical guidelines on labour.} https://www.health.qld.gov.au/
    \item NHS. \textit{Signs that labour has started.} https://www.nhs.uk/pregnancy/labour-and-birth/signs-of-labour/introducing-labour/
    \item Mayo Clinic. \textit{Stages of labor and contraction tracking.} https://www.mayoclinic.org/healthy-lifestyle/labor-and-delivery/in-depth/stages-of-labor/art-20046545
\end{itemize}
